%!TEX root=../../note.tex
\subsection{Proof By Contradiction}
\begin{definition}
    Given a statement $A$, the logical expression $A \wedge \paren{\neg A }$ is called a contradiction.
\end{definition}
A contradiction is false in every row of its truth table. To prove a statement $S$ by contradiction,
assume $\neg S$ and derive a contradiction; since a true assumption cannot lead to a false
statement, $\neg S$ must be false, so $S$ is true. For an implication $A \implies B$, we assume
$A$ together with $\neg B$.

\begin{proposition}
    $\forall n \in \Z$, if $n^2$ is even, then $n$ is even.
\end{proposition}
\begin{proof}
    Let $n$ be an arbitrary integer.

    Assume $n^2$ is even.

    Suppose, for contradiction, that $n$ is odd.

    Then,
    \begin{equation*}
        \paren{\exists k \in \Z } n = 2k + 1.
    \end{equation*}

    It implies
    \begin{equation*}
        n^2 = \paren{2k + 1 }^2 = 4k^2 + 4k + 1 = 2(2k^2 + 2k ) + 1,
    \end{equation*}

    but $2k^2 + 2k \in \Z$ (since $k \in \Z$). So we have actually shown that $n^2$ is odd.

    This contradicts the assumption that $n^2$ is even.

    Thus, the assumption that $n$ is odd must be false. Thus, $n$ is even.
\end{proof}

\subsubsection*{Exercises}
\begin{itemize}
    \item Prove there is no greatest integer.
    \item Prove $\forall a, b \in \Z, a \geq 2 \implies \paren{a \nmid b } \vee \paren{a \nmid \paren{b + 1}}$.
    \item Prove $\sqrt{2}$ is irrational.
\end{itemize}

\subsection{Proving Uniqueness}
To prove ``there is a unique $x \in S$ such that $P(x)$'', we prove two things:
\begin{itemize}
    \item \textbf{Existence:} there is at least one $x \in S$ such that $P(x)$ is true
    (for example, by exhibiting one);
    \item \textbf{Uniqueness:} assume that $P(x)$ and $P(y)$ are both true for some $x, y \in S$,
    and prove that $x = y$.
\end{itemize}
Uniqueness alone does not guarantee that such an $x$ exists, so both parts are needed.

\begin{proposition}
    $\forall a, b \in \Z$, if $a \neq 0$ and $a \mid b$, then there is a unique integer $k$ such that
    $b = ka$.
\end{proposition}
\begin{proof}
    Let $a, b$ be arbitrary in $\Z$. Assume $a \neq 0$ and $a \mid b$.

    \textbf{Existence.} Since $a \mid b$, by definition of divisibility,
    \begin{equation*}
        \paren{\exists k \in \Z } b = ka.
    \end{equation*}

    \textbf{Uniqueness.} Suppose $k, m \in \Z$ satisfy
    \begin{equation*}
        b = ka \qquad \text{and} \qquad b = ma.
    \end{equation*}

    Thus, $ka = ma$.
    Since $a \neq 0$, we may divide both sides by $a$, and we have
    \begin{equation*}
        k = m.
    \end{equation*}

    Thus, there is a unique $k \in \Z$ such that $b = ka$.
\end{proof}
\begin{remark}
    The hypothesis $a \neq 0$ is exactly what the uniqueness step uses. If $a = 0$ and $b = 0$,
    then $0 \mid 0$, but \emph{every} integer $k$ satisfies $0 = k \cdot 0$.
\end{remark}

\subsubsection*{Exercises}
Use a proof by contradiction to show that two non-parallel lines have a unique intersection point.
